\section{Limitations}
\label{sec:limitations}

The exact theory is linear-Gaussian. Model M1 with A4a dense signal carries the capture law, the frontier predictions, the impossibility results, and the dominance theorem; nothing proved here extends automatically to nonlinear linkages, discrete outcomes, or non-Gaussian factor structure. What supports broader scope today is empirical: a $28$-variant robustness panel covering $t$-distributed errors, Rademacher-half loadings, row-heteroskedastic noise, correlated factors, $r \pm 1$ misspecification, and sparse confounding leaves the qualitative ordering of estimators and the phase-diagram structure intact in every arm. The dedicated work package intended to convert these panels into a theoretical universality statement was cut before proof, so non-Gaussian protection should be read as conjecture backed by panels, not as theorem.

The benchmark negative controls are conservative from below. Realized control sizes ($0.000$ to $0.003$) sit beneath the predeclared $[0.02, 0.10]$ band's lower edge, mechanistically explained by permutation and split-half evaluation removing the beta-leakage component of the twin-null variances. The dangerous direction never fires (zero anti-conservative rejections across more than $3{,}000$ control evaluations), so the failure is of literalism rather than validity, but it does mean realized size on real geometry is below nominal and the reported power figures inherit that conservatism.

Frontier prediction is alignment-dependent. Concentrating the link on the dominant mega-spike direction collapses detectability at fixed strength (power $0.05$ versus $0.48$ for the spread placement), through the leakage-tax mechanism absorbed by that coordinate's twin scale; prediction quality is trustworthy for spread links and degrades under extreme concentration. Practitioners receive placement guidance from the package, not a uniform guarantee.

The computational probe is not an information-theoretic bound. Blindness claims are class-scoped by construction: purely covariate-based alarms are blind by data processing, and the shipped spike-coordinate alarm obeys the proved power ceiling of Proposition~\ref{prop:ceiling}, confirmed in nine of nine blind strata. Our stronger plan, an OMH-style mutual-contiguity statement for the full augmented spectrum, was falsified by its own pre-registered tests (Section~\ref{sec:blindness}), so no universal eigenvalue-impossibility is claimed anywhere; the gradient-boosting probe on a frozen feature map reads the response-scale signature at low aspect ($c = 0.2$, AUC up to $1.0$), which is precisely why the frozen universal-invisibility declaration failed its own gate and was rescoped. A machine-learning probe passing chance-level is evidence, not a lower bound.

Two structural caveats concern what the benchmarks can and cannot test. First, the injections build in exactly the loading geometry the theory assumes (dense Haar loadings on real covariate spectra), so the benchmarks validate calibration, frontier placement, and the adjustment recipe on realistic spectra; they cannot falsify the dense-loading assumption itself, and a practitioner facing sparse or discretely structured loadings is outside the validated envelope until robustness panels of that kind are run. Second, our alignment evidence sweeps the confounding direction's angle to the factor axes, and the benchmark sensitivity battery varies link concentration, but we ran no dedicated size sweep under deliberate correlation between $\beta$ and the top eigenvectors, the sharpest form of A4a violation; how fast alarm size inflates along that axis is an open measurement we flag rather than guess. Relatedly, the twin-and-permutation calibration inherits its generative model (Haar loadings, homoskedastic bulk edge, dropped factor link); only response-noise heteroskedasticity is stress-tested on this axis, so size guarantees are as strong as that model on a given dataset.

Finally, treatment-effect validation rests on a single economics family: the k401ksubs panel carries all M2 and PF-3 evidence on real data, corroborated by simulator cells but not yet by a second applied domain. And $r$ misspecification, while least damaging among the variants tested (median bias ratio $1.48$, injection-side powers diluted only moderately at $0.28$ to $0.44$), is tolerated within about one spike, not unbounded; the package therefore reports its rank estimate alongside every verdict rather than treating $r$ as known.

\section{Discussion}
\label{sec:discussion}

The results are orthogonal to sensitivity-analysis methodology rather than competitive with it. Omitted-variable sensitivity analyses in the style of Cinelli and Hazlett \citep{cinelli2020making} and E-value approaches \citep{vanderweele2017evalue} answer a conditional question: how strong would a hypothetical confounder have to be to overturn a conclusion. They quantify consequences given existence. This paper answers a complementary question: does the data itself contain spectral evidence that a dense confounder is present, above what strength would such evidence appear, and when does non-detection carry information rather than merely certify blindness. The combined workflow ran end to end on the benchmarks; screen with the calibrated alarm, adjust by hard trimming when it fires, and hand the residual hypothetical risk to sensitivity analysis when the blind-region certificate says non-detection means nothing. Neither tool subsumes the other, and the certificate tells the analyst which regime they are in, which is what a practitioner needs before trusting an adjusted estimate either way.

Three extensions seem most valuable. First, a non-Gaussian universality theorem replacing the empirical variant panels would convert the strongest conjectural claim into a proved one. Second, the fixed-$r$ spike analysis should extend to growing $r$, where the frontier algebra suggests qualitatively new scaling. Third, time-series spectra with serial dependence would test whether the twin-calibration idea survives structure beyond independent rows, the setting where batch-like spikes are most common in practice.
